


\documentclass[11pt]{article}
\usepackage{amsgen,amsmath,amstext,amsbsy,amsopn,amsfonts,amsthm,amscd}
%\usepackage{bbm}
\usepackage[dvips]{graphicx}
\usepackage{enumerate}
\setlength{\textwidth}{16.0cm} \setlength{\textheight}{25cm}
\voffset=-2.50cm \hoffset=-2.00cm \setlength{\topmargin}{0.4cm}

%%%%%%%%%%%%%%%%%%%%%%%%%%%%%%%%%%%%%%%%%%%%%%%%%%%%%%%%%%%%%%

\newtheorem{lem1}{Lema}%[chapter]
\newtheorem{th1}{Teorema}%[chapter]
\newtheorem{prop1}{Proposici\'{o}n}%[chapter]
\newtheorem{defi1}{Definici\'{o}n}%[chapter]
\newtheorem{coro1}{Corolario}%[chapter]
\newtheorem{ob1}{Observaci\'{o}n}%[chapter]
\newtheorem{obs1}{Observaci\'{o}n}%[chapter]
\newtheorem{nota1}{Nota}%[chapter]
\newtheorem{eje1}{Ejemplo}%[chapter]
\newtheorem{ejes1}{Ejemplo}%[chapter]
\newtheorem{ejerc1}{Ejercicio}%[chapter]
\newtheorem{ejercs1}{Ejercicio}%[chapter]
\newenvironment{prop}{\begin{prop1}\rm}{\end{prop1}}
\newenvironment{defi}{\begin{defi1}}{\end{defi1}}
\newenvironment{coro}{\begin{coro1}}{\end{coro1}}
\newenvironment{th2}{\begin{th1}}{\end{th1}}
\newenvironment{nota}{\begin{nota1}\rm}{\end{nota1}}
\newenvironment{lem}{\begin{lem1}}{\end{lem1}}
\newenvironment{ob}{\begin{ob1}\rm}{\end{ob1}}
\newenvironment{obs}{\begin{obs1}\rm}{\end{obs1}}
\newenvironment{eje}{\begin{eje1}\rm}{\end{eje1}}
\newenvironment{ejes}{\begin{ejes1}\rm}{\end{ejes1}}
\newenvironment{ejerc}{\begin{ejerc1}\rm}{\end{ejerc1}}
\newenvironment{ejercs}{\begin{ejercs1}\rm}{\end{ejercs1}}
%%%%%%%%%%%%%%%%%%%%%%%%%%%%%%%%%%%%%%%%%%%%%%%%%%%%%%
%\DeclareSymbolFont{cucu}{U}{bbmss}{m}{n}
\DeclareMathOperator{\CC}{\mathbb C}%{\mathalpha}{cucu}{67}
\DeclareMathOperator{\NN}{\mathbb N}%{\mathalpha}{cucu}{78}
\DeclareMathOperator{\RR}{\mathbb R}%{\mathalpha}{cucu}{82}
\DeclareMathOperator{\QQ}{\mathbb Q}%{\mathalpha}{cucu}{81}
\DeclareMathOperator{\ZZ}{\mathbb Z}%{\mathalpha}{cucu}{90}
%%%%%%%%%%%%%%%%%%%%%%%%%%%%%%%%%%%%%%%%%%%%%%%%%%%%%%
\DeclareMathOperator{\sen}{sen}
\DeclareMathOperator{\arcsen}{arcsen}
\DeclareMathOperator{\senh}{senh} \DeclareMathOperator{\fr}{Frac}
%%%%%%%%%%%%%%%%%%%%%%%%%%%%%%%%%%%%%%%%%%%%%%%%%%%%
\newcommand{\ran}[1]{{\em Im\/}(#1)}
\newcommand{\ca}[1]{{\em card\/}(#1)}
\newcommand{\sop}[1]{{\em supp\/}(#1)}
\newcommand{\tra}[1]{{\em traza\/}(#1)}
\newcommand{\dis}[1]{d (#1)}
\newcommand{\inte}[1]{{\em Int\/}(#1)}
\newcommand{\conv}[1]{{\em conv\/}(#1)}
\newcommand{\too}{\longrightarrow}
\newcommand{\oo}{\overline}
\newcommand{\cq}{\begin{flushright} $\Box$ \end{flushright}}
\newcommand{\cqd}{\, \, \rule{7pt}{7pt}}
\def\limsup{\mathop{\overline{\rm lim}}}
\def\liminf{\mathop{\underline{\rm lim}}}
\renewcommand{\theenumii}{\arabic{enumii}}
\renewcommand{\labelenumii}{\theenumi.\theenumii.}
\renewcommand{\theenumiii}{\arabic{enumiii}}
\renewcommand{\labelenumiii}{\theenumi.\theenumii.\theenumiii}
%%%%%%%%%%%%%%%%%%%%%%%%%%%%%%%%%%%%%%%%%%%%%%%%%%%%%%%%%%%%%%%%%
\begin{document}
\noindent 
\section{ ANÁLISIS DEL CASO DE DOS CIRCUNFERENCIAS SEGÚN EL ARTÍCULO DE ZEROS DEL ARCHIV} 

\noindent Resultados sobre la acaotación de $\mathal{D}$ en el 

\noindent {\bf CASO CONTINUO DOS CIRCUNFERENCIAS} : 


\begin{enumerate} 
\item cuando la circunferencia segunda está contenida en el disco abierto primero referencia artículo de Archiv ejemplo 1: 
CASO SECUENCIALMENTE DOMINADO.....Y POR TANTO LA D ESTÁ ACOTADA. 
\item cuando el centro de la circunferencia segundo está contenido en el disco de la primera...puede tener radio muy grande...D está acotada (referncia artículo de Archiv ) NO ES UN CASO SECUENCIALMENTE DOMINADO NI MATRICIALMENTE SECUENCIALMENTE 
\item Cuando el centro de la segunda no está contenido en el disco de la primera entonces no sabemos qué pasa con D ????NO SECUENCIALEMENTE DOMINADO. 
\end{enumerate} 

\noindent {\bf CASO DE CIRCUNFERENCIA Y ATOMO.} 
\begin{enumerate} 
\item cuando el átomo está dentro del disco abierto  la D está acotada y 
\begin{enumerate} 
\item Puesto que $a$ pertenece al disco es un caso SECUENCIAlMENTE DOMINADA...COROLARIO 36
\item cuando el átomo está fuera del disco abierto...D no sabemos qué pasa....pero parece no acotado en la experimentación. 
\end{enumerate} 






\section{CUANDO INTRODUCES LA MATRIZ D} 
\noindent Consideremos el operador multiplicaci\'{o}n por $z$ 
$$
\mathcal{D}: \mathbb{P}[t] \to \mathbb{P}[z], \quad \mathcal{D}(p(t))=zp(z).
$$

\noindent La representación de este operador en la base de polinomios ortogonales de Sobolev $\{P_n(z)\}_{n=0}^{\infy}$ es una matriz de Hessemberg superior ya que 
$zP_n(z)$ es un polinomio de grado $n+1$ que se puede expresar en términos de los polinomios $\{P_0(z),P_1(z), \dots,P_n(z),P_{n+1}(z)\}$ de la forma: 
$$
zP_n(z)= d_{0,0}P_0(z)+d_{1,0}P_1(z)+\dots+d_{n,0}P_n(z)+d_{n,n+1}P_{n+1}(z) 
$$

\noindent Consideremos la matriz infinita $\mathbf{D}=(d_{i,j})_{i,j=0}^{\infty}$ y llamaremos $\mathbf{D}_n$ a la sección de orden $(n+1)\times (n+1)$. 
\noindent Probaremos que: 
\begin{enumerate} 
\item los ceros de $P_{n+1}(z)$ son  autovalores de $\mathbf{D}_n$. 
\item El determinante de $\mathbf{D}_n -z \mathbf{I}_n$ denotado por $\vert \mathbf{D}_n -z \mathbf{I}_n\vert $ donde $\mathbf{I}_n$ es la matriz identidad de dimensiones $(n+1)\times (n+1)$ es un polinomio proporcional a $P_{n+1}(z)$, es decir, tienen las misma raíces.
\end{enumerate} 


\noindent PRUEBA DE QUE SI $Z_0$ es raíz de $\varphi_n(z)$ entonces $z_0$ es autovalor de $\mathbf{D}_n^{t}$ y por tanto de $\mathbf{D}_n$. 

\begin{proof} 
\noindent Sea $z_0$ una raíz de $P_{n+1}(z)$, es decir, $P_{n+1}(z_0)=0$. Entonces: 

$$
z_0P_0(z_0)= d_{0,0}P_0(z_0)+d_{1,0}P_1(z_0)
$$
$$
z_0P_1(z_0)= d_{0,1}P_0(z_0)+d_{1,1}P_1(z_0)+d_{2,1)P_{2}(z_0)
$$
$$ 
----------------
$$
$$
z_0P_n(z_0)= d_{0,n}P_0(z_0)+d_{1,n}P_1(z_0)+\dots+d_{n,n}P_n(z_0)+d_{n+1,n}P_{n+1}(z_0) 
$$

\noindent Puesto que $P_{n+1}(z_0)=0$, la última igualdad queda: 
$$
z_0P_n(z_0)= d_{0,0}P_0(z_0)+d_{1,0}P_1(z_0)+\dots+d_{n,0}P_n(z_0). 
$$

\noindent Esto se puede expresar matricialmente de la siguiente forma 
$$
z_0\begin{pmatrix} P_0(z_0) \\ P_1(z_0) \\  \vdots \\ P_n(z_0) \end{pmatrix} = \mathbf{D}^{t}_n \begin{pmatrix} P_0(z_0) \\ P_1(z_0) \\  \vdots \\ P_n(z_0) \end{pmatrix}
$$

\noindent lo cual significa que $z_0$ es autovalor de la matriz $\mathbf{D}^{t}_n$ siendo el autovector $(P_0(z_0),P_1(z_0),\dots,P_n(z_0))$ (que es no nulo ya que $P_0(z)$ es constante). 

\noindent Por lo tanto, todos las raíces de $P_n(z_0)$ son autovalores de la matriz $\mathbf{D}^{t}_n$. 

\noindent 
\end{proof} 

\noindent Como consecuencia El determinante de $\mathbf{D}_n -z \mathbf{I}_n$ denotado por $\vert \mathbf{D}_n -z \mathbf{I}_n\vert $ donde $\mathbf{I}_n$ es la matriz identidad de dimensiones $(n+1)\times (n+1)$ es EL POLINOMIO MÓNICO.  



\section{COMENTARIOS ADICCIONALES} 
\begin{enumerate} 
\item El análisis de casos refleja  y constata los resultados sobre la acotación de $D$ (en el caso de dos circunferencias y átomos) de nuestros artículos....sobre la acotación del operador D en los espacios de Sobolev.
\item Resaltar la convergencia de los $\rho(D_n)$, que es novedoso, pues en el caso de conocerse que los ceros están acotados  no tendría por qué ocurrir que este límite exista y esto ha queddado reflejado EN TODOS LOS EJEMPLOS.  


\noindent Otro problema sería el de la convergencia de los números complejos donde se alcanzan estos  valores máximos ....que serían 
por ejemplo los llamas $z_n$ siendo la definición de este: 


$z_n\in Z_n$ tal que $\vert z_n \vert = \rho(\mathbf{D}_n)$.  


que yo no me acuerdo de lo que ocurría. 

\item Resaltar que es la primera vez que se  ha estudiado el problema de la convergencia del {\it segundo} valor singular de las matrices $D_n$ en el contexto de la localización o acotación de  ceros de polinomios Sobolev, o en relación con los soportes de las medidas involucradas.  

\noindent Cuando pongas la conjetura de que en el caso de $\Vert \mathcal{D}\Vert = \infty$ y 
$$
\lim_{n\to \infty} \sigma_{2,n}=R
$$

\noindent {\bf es importante añadir que}: 

\noindent  en el caso standard de polinomios ortogonales asociados a una única medida con soporte compacto, es decir, en el caso de 
$$
<p,q>= \int \vert p(z) \vert^{2} \; d\mu_1
$$


 \noindent  es conocido que: 
 $$
 \lim_{n\to \infty} \sigma_{1,n} = \lim_{n\to \infty} \Vert \mathbf{D}_n \Vert = R 
 $$
 \noindent siendo $R= \max\{ \vert z \vert : z\in {\it supp}(\mu)\}$.  

\end{enumerate} 
 
\end{document} 

\noindent Producto Sobolev continuo: caso secuencialmente dominado. 

\noindent Estimaci\'{o}n de la cota de acotaci\'{o}n de ceros de polinomios ortogonales de Sobolev. 

\noindent Consideremos el siguiente producto de Sobolev: 
$$
<p(t),q(t)>=  \int_{a}^{b} p(t)q(t)w_0(t)\; dt +  \int_{a}^{b} p'(t)q'(t)w_1(t)\; dt
$$

\noindent donde $w_0(t),w_1(t)$ son pesos continuos de modo que est\'{a}n sequencialmente dominados en el sentido de que: Existe una constante $C>0$ tal que: 
$$
w_1(t) \leq Cw_0(t) , \qquad t\in [a,b]
$$

\noindent Puesto que $\dfrac{w_1(t)}{w_0(t)}$ es una funci\'{o}n continua en $[a,b]$ consideremos 
$$
\alpha= \max \{ \dfrac{w_1(t)}{w_0(t)}, t\in [a,b]\}.
$$

\noindent Consideremos el operador multiplicaci\'{o}n por $t$ 
$$
\mathcal{D}: \mathbb{P}[t} \to \mathbb{P}[t], \quad \mathcal{D}(p(t))=tp(t).
$$

\noindent Veamos que $\mathcal{D}$ es acotado en este caso adem\'{a}s: 
$$
 \Vert \mathcal{D} \Vert \leq C= (R^{2}+2\alpha)^{1/2}.
 $$

$$
\Vert \mathcal{D}(p(t))\Vert^{2} = \int_{a}^{b} t^2p^2(t)w_0(t)\; dt +  \int_{a}^{b} (tp'(t))^2w_1(t)\; dt= \int_{a}^{b} t^2p^2(t)w_0(t)\; dt +  \int_{a}^{b} (p(t)+tp'(t))^2w_1(t)\; dt \leq 
$$
\noindent utilizando que $(a+b)^2\leq 2(a^2+b^2)$ si $a,b\geq 0$ se tiene que: 
$$
 \int_{a}^{b} t^2p^2(t)w_0(t)\; dt+2\int_{a}^{b} p(t)^2w_1(t)\; dt +\int_{a}^{b} +t^2p'(t))^2w_1(t)\; dt \leq
 $$
 $$
  R^{2}\left( \int_{a}^{b} p^2(t)w_0(t)\; dt +  \int_{a}^{b} (p'(t)^2w_1(t)\; dt+ 2\int_{a}^{b} p(t)^2w_1(t)\; dt \leq 
$$
$$
R^{2} \Vert p(t)\Vert^{2} + 2\alpha \int_{a}^{b} p(t)^2w_0(t)\; dt \leq (R^{2}+2\alpha)\Vert p(t) \Vert^{2}
$$

\noindent Por ejemplo, si $[a,b]=[-1,1]$ y $\alpha =1$ quedar\'{\i}a $C=\sqrt{1+2}=\sqrt{3}$. 

\begin{enumerate} 
\item Experimentar con ejemplos sequencialmente dominados $w_0(t)=t^n$ y $w_1(t)=t^{2n}$ con $n$ par en el intervalo $[-1,1]$, se aprecia que la c ota es buena ???
\item Experimentar con ejemplos que no sean sequencialmente dominados $w_0(t)=t^{2n}$ y $w_1(t)=t^{n}$ con $n$ par en el intervalo $[-1,1]$, se aprecian resultados an\'{a}logos al caso anterior. 
\item Observar que en los casos anteriores todos los ceros están en el círculo cerrado de centro el origen y radio $\sqrt{3}$. 
\end{enumerate} 

\section{CASO DE CIRCUNFERENCIAS} 




\noindent Consideramos las medidas ${\bf m}_{a;r}$ que son las medidas de Lebesgue normalizadas en las circunferencias de centro $a$ y radio $r$ por lo que: 
$$
<p(z),q(z)>_{{\bf m_{a;r}}} = \int_{0}^{2\pi}\dfrac{1}{2\pi} p(a+re^{it})\overline{q(a+re^{it})} \; dt
$$

\noindent Vamos a estudiar el comportamiento asintótico de autovalores, valores singulares, etc de productos de Sobolev 

\noindent Utilizando que los ceros se transforman por semejanzas y homotecias,(Carmen) podemos estudiar los casos más sencillos: 
\begin{enumerate} 
\item Productos de Sobolev continuos asociados a medidas de Lebesgue. 

$$
<p,q>_{\mathbf{S}} =\int p(z) \overline{q(z)}d{\bf m}_{0;r} + \int_{\mathbf{S}_r} p'(z) \overline{q'(z)}d{\bf m_{a;s}}
$$

\noindent estudiando los distintos casos dependiendo de la posición de las circunferencias. 
\item Productos de Sobolev discretos con un átomo. 
$$
<p,q>=\int p(z) \overline{q(z)}d{\bf m}_{a;r} + p(0)\overline{q'(0)}.
$$

\noindent dependiendo de la posición del átomo con respecto a la circunferencia. 
\end{enumerate} 


\noindent En todos los casos  se trata de estudiar: 
\begin{enumerate} 
\item Si existe y cuanto es 
$$
\lim_{n\to \infty} \rho(\mathbf{D_n})
$$

\item 
$$
\lim_{n\to \infty} \sigma_{1,n} = \Vert \mathcal{D} \Vert 
$$

\item en el caso de $\Vert \mathcal{D} \Vert = \infty$, 
$$
\lim_{n\to \infty} \sigma_{2,n} = R??????? 
$$
\end{enumerate} 






\section{Resultados conocidos en el artículo } 

\noindent Utilizando el artículo 

\noindent $\bullet$ Escribano, Gonzalo, {\it Zeros of Sobolev polynomials and some matrix inequalities, arXiv buscar referencia}

\noindent {\bf Resultados conocidos para el caso continuo.} 
$$
<p,q>_{\mathbf{S}} =\int p(z) \overline{q(z)}d{\bf m}_{a;r} + \int_{\mathbf{S}_r} p'(z) \overline{q'(z)}d{\bf m_{b;s}}
$$
\begin{enumerate} 

\item En el {\bf Ejemplo 1} se incluye la situación de Sobolev continuo con ${\bf m}_{b;s} \subset {\bf m}_{a;r}$ que, es un caso MATRICIAL secuencialmente continuo. 
\item Utilizando la {\bf Proposición 6}, en el caso Sobolev continuo con EL CENTRO DE LA CIRCUNFERENCIA SEGUNDA EN EL INTERIOR DE LA CIRCUNFERNCIA PRIMERA $b\in \mathbb{D}(a;r)$ y radio $s$ CUALESQUIERA  se prueba que, aunque no es un caso secuencialmente continuo, sin embargo el operador multiplicación es acotado. 
\item Quedarían por estudiar los casos en los que EL CENTRO DE LA SEGUNDA NO ESTÁ EN EL INTERIOR DE LA PRIMERA $b\notin \mathbb{D}(a;r)$: disjuntas, tangentes, o que se corten. De especial interés son los casos en los que se consideran la medida de Lebesgue y la medida centrada en $1$ y de radio $1$ (circunferencias tangentes). 


\end{enumerate} 

\noindent {\bf Resultados conocidos para el caso discreto} 

$$
<p,q>=\int p(z) \overline{q(z)}d{\bf m}_{a;r} + p(0)\overline{q'(0)}.
$$

\begin{enumerate} 
\item En el caso del átomo EN EL INTERIOR DE LA CIRCUNFERENCIA, QUE ES UN CASO SECUENCIALMENTE DOMINADO, es decir, $0\in \mathbb{D}(a;r)$ se tiene que el operador multiplicación es acotado. 

\item EN EL CASO DE QUE EL CENTRO NO ESTÁ EN EL INTERIOR DE LA CIRCUNFERENCIA YO CREO, AUNQUE NO ESTÁ PROBADO que si el átomo está fuera del disco abierto entonces el operador multiplicación no está acotado y es lo que se ve en los ejemplos. 

\end{enumerate} 


\end{document} 
\section{Caso discreto con átomo en $0$.} 

\noindent Cuando el átomo es el cero, tenemos unos resultados interesantes: 

\noindent Consideremos el subespacio de los polinomios $\mathbb{P}_0 \cap \mathbb{P}'_0={\ker}\delta_{0} \cap {\ker}\delta^{(1)}_0$ que es un subespacio de codimensión 2, se tiene que: 

\noindent Este subespacio es precisamente 
$$
\mathbb{P}_0 \cap \mathbb{P}'_0=\{z^2p(z) : p(z)\in \mathbb{P}[z]\}.
$$

\noindent En efecto, los polinomios en $\mathbb{P}_0 ={\ker}\delta_{0}$ de grado $n\geq 2$ son polinomios del tipo 
$$
a_2z^2+a_3z^3+\dots+a_nz^{n}= z^2(a_2+\dots+a_{n}z^{n-2}). 
$$

\noindent Si miramos las normas Sobolev en $\mathbb{P}_0 \cap \mathbb{P}'_0$: 

$$
\Vert \mathcal{D}(p(z))\Vert^{2}_{\mathbf{S}}= \int \vert zp(z)\vert^{2}\; d\mu + \vert \delta^{(1)}_{0}(zp(z))\vert^{2} 
= \Vert zp(z) \Vert^{2}_{\mu} + \delta_{0}(p(z)+zp'(z))= \Vert zp(z) \Vert^{2}_{\mu} 
$$

\noindent y por otra parte,
$$
\Vert p(z) \Vert^{2}_{\mathbf{S}} = \Vert p(z)\Vert^{2}_{\mu} + p'(0)^{2}= \Vert p(z) \Vert^{2}_{\mu}
$$

\noindent por lo tanto:
\begin{prop} Sea $\mu$ una medida con  soporte acotado, es decir, $\mathcal{D}$ es acotado con la norma $\Vert \cdot \Vert_{\mu}$. Consideremos la norma Sobolev 
$$
\Vert \cdot \Vert^2_{\mathbf{S}} = \int \vert p(z) \vert^{2} \; d\mu + \vert p'(0) \vert^{2}
$$
Consideramos el espacio vectorial 
$$
X_0=\mathbb{P}_0[z]\cap \mathbb{P}'_0[z]= \{p(z):\; p(0)=0\}= \{zp(z):\; p(z)\in \mathbb{P}[z]\}.
$$
\noindent Entonces 
\begin{enumerate}
\item Sea $X_0=\mathbb{P}_0[z]\cap \mathbb{P}'_0[z]$, los espacios  $(X_0,\Vert \cdot \Vert_{\mu})$ y $(X_0,\Vert \cdot \Vert_{\mathbf{S}})$ son el mismo por lo que: 
$$\Vert \mathcal{D} \vert X_0 \Vert_{\mathbf{S}}  =  \Vert \mathcal{D}\Vert_{\mu}.$$

\item $\mathcal{D}$ acotado en $(\mathbb{P}_0[z], \Vert \cdot \Vert_{\mathbf{S}})$ siendo 
$$
\Vert \mathcal{D}\vert \mathbb{P}_0[z] \Vert_{\mathbf{S}} = \Vert \mathcal{D}\Vert_{\mu}.
$$


\end{enumerate} 
\end{prop} 

\begin{proof} 






\noindent Si $p(z)\in X_0$ se tiene que: 

$$
\Vert \mathcal{D}(p(z))\Vert^{2}_{\mathbf{S}}= \int \vert zp(z)\vert^{2}\; d\mu + \vert \delta^{(1)}_{0}(zp(z))\vert^{2} 
= \Vert zp(z) \Vert^{2}_{\mu} + \delta_{0}(p(z)+zp'(z))= \Vert \mathcal{D}(p(z)) \Vert^{2}_{\mu} 
$$

\noindent y por otra parte,
$$
\Vert p(z) \Vert^{2}_{\mathbf{S}} = \Vert p(z)\Vert^{2}_{\mu} + \vert p'(0)\vert^{2}= \Vert p(z) \Vert^{2}_{\mu}
$$







\noindent Luego, 
$$
\sup \{ \Vert \mathcal{D}\vert X_0(p(z)\Vert_{\mathbf{S}}: \; \Vert p(z)\Vert_{\mathbf{S}}=1\} = \sup \{ \Vert \mathcal{D}\vert X_0(p(z)\Vert_{\mu}: \; \Vert p(z)\Vert_{\mu}=1\} 
$$

\noindent Es decir, 
$$
\Vert \mathcal{D}\vert X_0\Vert_{\mathbf{S}}= \Vert \mathcal{D}\vert X_0\Vert_{\mu}
$$

\noindent Por lo tanto: 
$$
\Vert \mathcal{D}\vert X_0\Vert_{\mathbf{S}}\leq \Vert \mathcal{D}\Vert_{\mu}
$$





\noindent Ahora bien, es conocido que para el caso de medidas 
$$
\Vert \mathcal{D} \Vert_{\mu} =\lim_{n\to \infty} \Vert \mathcal{D}(\dfrac{z^n}{\Vert z^{n}\Vert_{\mu}})\Vert_{\mu} 
$$

\noindent eso significa que puesto que los $z^n\in X_0$, se tiene que $\Vert z^n\Vert_{\mathbf{S}}= \Vert z^{n}\Vert_{\mu}\Vert$ por lo que: 
$$
\Vert \mathcal{D}\vert X_0 \Vert_{\mathbf{S}} \geq \Vert \mathcal{D}(\dfrac{z^n}{\Vert z^{n}\Vert_{\mathbf{S}}})\Vert_{\mathbf{S}}= \mathcal{D}(\dfrac{z^n}{\Vert z^{n}\Vert_{\mu}})\Vert_{\mu} 
 $$

\noindent y pasando al límite: 
$$
\Vert \mathcal{D}\vert X_0 \Vert_{\mathbf{S}}= \Vert \mathcal{D}\vert_{\mu} 
$$
\noindent por lo que 
$$
\Vert \mathcal{D}\vert X_0\Vert_{\mathbf{S} } = \Vert \mathcal{D}\Vert_{\mu}
$$ 


\end{proof} 








\noindent Si consideramos el segundo número de Gelfand, 

$$
c_2(\mathcal{D})=\inf \Vert \mathcal{D}\vert X\Vert : X \subset \mathbb{P}[z], cod(X)=1\} \leq \Vert \mathcal{D}\vert \mathbb{P}_0[z]\Vert = R.
$$

\noindent se tiene que si $\mu$ tiene soporte acotado, entonces 
$$
c_2(\mathcal{D}) \leq R.
$$

\noindent {\bf PROBLEMAS QUE SURGEN:} 
\noindent Se puede probar que siempre se tiene que $\sigma_{2,n}\leq c_{2}(\mathcal{D})$ para todo $n\in \NN$ por lo que 
$$
\lim_{n \to \infty} \sigma_{2,n}\leq  c_{2}(\mathcal{D})
$$
\begin{enumerate} 
\item Es cierto que $c_2(\mathcal{D})=R$?????? 
\item Es cierto que 
$$
\lim_{n\to \infty} \sigma_{2,n}=c_2(\mathcal{D}).?????????
$$

\end{enumerate} 



\begin{prop} Consideremos el producto de Sobolev discreto: 
$$
\Vert p(z) \Vert^{2}_{\mathbf{S}}= \int \vert p(z) \vert^{2} \; d\mu + \vert p'(0) \vert^{2}.
$$

\noindent Son equivalentes: 
\begin{enumerate} 
\item $\mathcal{D}$ es acotado con $\Vert \cdot \Vert_{\mathbf{S}}$, es decir, 
$\Vert \mathcal{D} \Vert_{\mathbf{S}}< \infty.$
\item $0$ es un punto de evaluación continua de $\mu$, es decir,  existe $C>0$ tal que 
$$
\vert p(0) \vert^{2} \leq C \int \vert p(z) \vert^{2} d\mu. 
$$

\end{enumerate} 


\begin{proof} Si $0$  es un punto de evaluación de la medida $\mu$, es decir,   existe $C>0$ tal que
$$
\vert p(0) \vert^{2} \leq C\int \vert p(z) \vert^{2}d\mu, \qquad p(z) \in \mathbb{P}[z].
$$

\noindent Entonces por los resultados de EG1 se tiene que $\{\mu,\delta_{0}\}$ es secuencialmente dominado por lo que $\mathcal{D}$ es acotado con la norma $\Vert \cdot \Vert_{\mathbf{S}}$.

\noindent Supogamos ahora que $\mathcal{D}$ es acotado pero que  $0$ no es punto de evaluación continuo de la medida.  Puesto que $\mathcal{D}$ es acotado por los resultados de EG1 se tiene que 
$$
\vert p(0) \vert^{2} \leq C \Vert p(z) \Vert^{2}_{\mathbf{S}} = C( \int  \vert p(z) \vert^{2}d\;\mu + \vert p'(0)\vert^{2}.
$$

\noindent Entonces si consideramos 
$$
\mathbb{P}'_0[z]=\overline{\{p(z): p'(0)=0\}}
$$
\noindent se trata de un espacio de codimensión menor o igual que $1$. pero una aplicación lineal definida en un subespacio de codimensión $1$ en $\mathbb{P}[z]$ puede no ser acotada.!!!!!!!!

\end{proof} 
\noindent El resultado anterior es consecuencia del siguiente resultado: 
\begin{prop}  Sea $L:\mathcal{H}\to \CC$ una aplicación lineal en un espacio de Hilbert $\mathcal{H}$. Si existe un subespacio de codimensión finita $V$ tal que $L\vert V$ es continua, entonces $L$ es continua en $\mathcal{H}$. 
\end{prop} 
\begin{proof} Todo subespacio de codimensión finita está complementado en un espacio de Hilbert. De hecho, 
$$
\mathcal{H}= V \bigoplus V^{\perp} 
$$

\noindent Por lo tanto, si $u\in \mathcal{H}$ se tiene que $u=P_{V}(u)+P_{V^{\perp}}(u)$. Puesto que $V$ es finito dimensional, la aplicación lineal es acotada en $V$, es decir,  $\Vert L \vert V \Vert<\infty$. Entonces  considerando $M= \max \{ \Vert L\vert V \Vert, \Vert L\vert V^{\perp}\Vert \}$ se tiene que, puesto que las proyecciones ortogonales tienen norma menor o igual que $1$, 
$$
\Vert L(u)\Vert = \Vert L(P_{V}(u))+L(P_{V^{\perp}}(u))\Vert \leq 
$$
$$
\Vert L\vert V (P_{V}(u)\Vert + \Vert L\vert V (P_{V^{\perp}}(u)\Vert +
2M \Vert u \Vert.
$$
\end{proof} 





\section{Hay aplicaciones lineales en espacios prehilbertianos acotadas en subespacios de codimensión finita pero no acotadas en todo el espacio. 
} 
\noindent Consideremos la siguiente matriz infinita, incluso de Hessemberg: 
$$
\mathbf{D} = \begin{pmatrix}
 1           &    1         & 1   & 1 & \hdots  \\
 \frac{1}{2} & \frac{1}{2}  & 0   & 0  & \hdots 
 \\ 
 0 & \frac{1}{3} & \frac{1}{3}   & 0  & \hdots 
\\ \vdots & \vdots & \vdots   & \vdots  & \ddots
\end{pmatrix} 
$$

\noindent $\mathbf{D}$ es una matriz de Hessemberg infinita, está definida en $c_{00}$ y claramente no define una aplicación lineal acotada en $c_{00}$ puesto que en el vector $u_n=(\frac{1}{\sqrt{n}},\dots,\frac{1}{\sqrt{n}},0,0,\dots)$ tiene norma una y sin embargo 
$$
\Vert \mathcal{D}(u_n) \Vert_{2}\geq \frac{n}{\sqrt{n}}\to \infty
$$

\noindent Sin embargo, si consideramos esta aplicación lineal en el subespacio vectorial de codimensión $1$, $[e_n,n\geq 1]$ se tiene que su representación es la matriz 

$$
\mathbf{D} = \begin{pmatrix}
 \frac{1}{2} & \frac{1}{2}  & 0   & 0  & \hdots 
 \\ 
 0 & \frac{1}{3} & \frac{1}{3}   & 0  & \hdots 
\\ \vdots & \vdots & \vdots   & \vdots  & \ddots
\end{pmatrix} 
$$

\noindent que define incluso un operador compacto !!!!!

\noindent Cómo cambia todo si el espacio es de Hilbert. 

\section{Una aplicación lineal EN UN HILBERT que está acotada en un subespacio vectorial de codimensión $1$ está acotada.} 

\noindent Poner en funcional, como existencia de las proyecciones ortogonales: 
si una aplicación lineal en un espacio de Hilbert es continua en un subespacio vectorial de codimensión finita entonces es continua en todo el espacio. 

\noindent Sea $\mathcal{D}$ una aplicación lineal DEFINIDA en un espacio de Hilbert. Supongamos que $\mathcal{D}$ está acotada en un subespacio cerrado $V$ de codimensión $1$. Entonces $\mathcal{D}$ está acotada en $\mathcal{H}$. 

\noindent Puesto que $V$ es un subespacio cerrado de dimensión $1$ se tiene que
$$
\mathcal{H}= V \bigoplus V^{\perp}
$$

\noindent Consideremos las proyecciones ortogonales $P_{V}$ y $P_{\V^{\perp}}$ las proyecciones ortogonales (	que existen por las propiedades de los espacios de HIlbert). 

\noindent Sea $\Vert \mathcal{D}/V\Vert =C$ entonces, para todo $u\in \mathcal{H}$ se tiene que 
$$
\Vert \mathcal{D}(u)\Vert =\Vert \mathcal{D}(P_V(u)+P_{V^{\perp}}(u)\Vert \leq \Vert \mathcal{D}(P_V(u)\Vert +\Vert \mathcal{D}(P_{V^{\perp}}(u))\Vert \leq 
$$
$$
\Vert \matcal{D}\vert V\Vert \Vert P_{V}(u)\Vert + \Vert \matcal{D}\vert V^{\perp}\Vert \Vert P_{V^{\perp}}(u)\Vert \leq \max\{\Vert \matcal{D}\vert V\Vert , \Vert \matcal{D}\vert V^{\perp}\Vert\}\Vert u \Vert 
$$

\noindent CLAVE: EXISTENCIA DE PROYECCIONES ORTOGONALES. 


\end{document} 
